% !TeX root = ../BTechProjectTemplate.tex
\chapter{Ethical Considerations and Social Impact}

\section{Introduction}
The technical evaluation of a model, using PSNR, SSIM, and LPIHS measures, only scratches the surface when the model is meant for clinical purposes. Even if a super-resolution model performs extremely well on all benchmarking standards, yet produces anatomically unfeasible results or cannot be used by the hospital centers where it could have the greatest impact, it is ultimately considered a failure. This chapter focuses on the ethical considerations associated with the Swin-MMHCA system, specifically its safety considerations, privacy considerations, and social impact.

\section{Clinical Safety and the Risk of Hallucination}
The first ethical consideration in relation to the generation of higher resolution images using medical imaging technology is that an artificial neural network may generate some details that seem realistic but are not representative of the patient's body. The issue is inherent to the nature of any generative process and relates to the likelihood of the generated results having little or nothing to do with reality.

The architecture of Swin-MMHCA was designed with the above issue in mind. There are several approaches used in order to make sure that the image generated by the network corresponds to the real image captured by a medical imaging device:

\begin{itemize}
    \item \textbf{Residual learning:} Instead of predicting the entire final image, the network is taught to make a prediction for the difference between the low-resolution input and the high-resolution desired result. This indicates that the low frequency structural information, which is the overall structure of the scan, is passed down from the original data and cannot be freely generated by the network.
    \item \textbf{Edge guidance:} The Sobel edge path is an explicit structure that comes from the input modes. The transformer has to honor these limits, which places a limit on where additional high-frequency information can go.
    \item \textbf{Multi-task semantic supervision:} The segmentation head and the lesion detector act as a mechanism of consistency. If there is a hallucination present in the upscaled image but not in the predicted segmentation map, then it will create a loss for the model to be consistent. Gradually, this helps the model learn not to hallucinate inconsistent structures.
\end{itemize}

Even with all these precautions, however, it may still be too soon to consider super-resolution imaging as the ground truth. When clinicians utilize such an approach, they must realize that the output will provide a magnified version of the scan and will help pinpoint areas of interest, but should never be used as the sole basis.

\section{Data Privacy and Patient Confidentiality}
The experiments performed in this study make use of the IXI database, which has undergone formal de-identification and public release. As such, the privacy concerns related to the IXI database are relatively simpler compared to the clinical application of the technology.
But implementing it in practice will not be an easy process. First of all, since three-dimensional MRI scans of the brain are essentially biometric information, because the surface of the brain can be recognized even after the face is removed using typical methods of de-identification. Thus, for example, any implementation of the algorithm in hospitals would involve advanced methods of de-identification of three-dimensional images and strict access restrictions.

One of the possible solutions to the problem of privacy and security concerns the application of federated learning techniques. The idea is that the model is not trained centrally, but independently at each institution and then aggregated using their gradients or weight updates. That way, Swin-MMHCA could be trained on much more diverse and extensive data than a single dataset allows.

\section{Healthcare Accessibility and Social Relevance}
There is a considerable distance between the theoretical capabilities of MRIs and practical accessibility across the world. Amongst the very few ways to bridge this divide without having to buy new hardware equipment is the use of software-driven super-resolution techniques.

In areas with limited resources, physicians have to rely on old 1.5T or low-field machines. The output from these devices is functional but does not match the resolution generated by high-field devices. Once it has been trained, Swin-MMHCA does not need any additional hardware to operate. It can run on a regular medical computer or, hypothetically, even on an online server connected via the internet. Thus, the advantages of high-field resolution would become available to organizations that would never afford a 3T device.

In addition, there are patient-related benefits that are sometimes overlooked in technical discourse. Prolonged scanning time is problematic not only in terms of convenience but because it actually constitutes a barrier to a considerable number of patients. Children with an inability to remain still, adults suffering from conditions like Parkinson's disease, and individuals with claustrophobia all encounter problems with prolonged imaging procedures. The use of technology capable of generating medically relevant images based on shorter and lower-resolution acquisitions leads to a decrease in the use of sedatives among children and a more pleasant experience for patients experiencing discomfort during long procedures.

From the standpoint of health care systems, shorter scanning periods lead to better performance, with lower costs and shorter waiting periods. This is important in areas where resources are scarce.

\section{The Human-AI Collaborative Model}
It should be made clear from the start what Swin-MMHCA can and cannot do. First and foremost, Swin-MMHCA is a preprocessing tool, one that supplies a sharper and cleaner image than the scanner alone. The task of making sense of it falls on the doctor, who still needs to perform all of the interpretive tasks manually. There is nothing in Swin-MMHCA's architecture that enables it to make a diagnosis, as this would run counter to its intended function.

The advantage of approaching it in such terms is an alleviation of the cognitive burden imposed on the radiologist. Poor-quality images require the interpretation process itself to correct for imperfections in the acquired data — interpreting the fuzzy features, ignoring noise and other artifacts, etc. With cleaner data, however, the doctor can devote more brainpower to what he was trained to do: identify uncommon patterns.

\subsection{Standardization and Regulatory Considerations}
Any AI-driven imaging technology will need to undergo assessment using criteria that surpass those defined in benchmarking databases prior to being deployed into clinical practice. A model could perform well in terms of PSNR on the IXI database but would nevertheless generate output images that can mislead clinical practitioners with regard to edge cases involving either rare pathologies or unusual anatomy or scan quality beyond what was observed during training.

Our intuition suggests that a pathological consistency metric is necessary as part of the evaluation criteria for a medical SR model. Pathological consistency refers to whether any medically significant information contained in the original image will persist or be altered in any way within the super-resolved output image.

It will become necessary for us to incorporate such an evaluation process as part of our regulatory review process, whether in terms of FDA processes in the United States or equivalent organizations in other countries, in order to build the trust necessary for clinical implementation.